\documentclass[11pt]{article}

\usepackage[T1]{fontenc}
\usepackage{amsmath,amssymb,amsthm}
\usepackage{booktabs}
\usepackage{geometry}
\usepackage{hyperref}
\usepackage{xurl}

\geometry{margin=1in}
\hypersetup{
  colorlinks=true,
  linkcolor=blue,
  citecolor=blue,
  urlcolor=blue
}

\newtheorem{theorem}{Theorem}[section]
\newtheorem{lemma}[theorem]{Lemma}
\newtheorem{proposition}[theorem]{Proposition}
\newtheorem{corollary}[theorem]{Corollary}
\theoremstyle{definition}
\newtheorem{definition}[theorem]{Definition}
\theoremstyle{remark}
\newtheorem{remark}[theorem]{Remark}

\newcommand{\gaminf}{\gamma^{\infty}}
\newcommand{\E}{E}
\newcommand{\V}{V}
\newcommand{\CNF}{\mathsf{CNF}}
\newcommand{\sha}{\operatorname{SHA256}}

\title{Certified Exclusion of Order-Twelve, Parameter-Three\\
Counterexamples to the \texorpdfstring{$\gamma$--$\theta$}{gamma--theta}
Conjecture}
\author{Author metadata to be supplied before submission}
\date{July 25, 2026}

\begin{document}

\maketitle

\begin{abstract}
In the one-guard-moves eternal domination game, an attack is made at an
unoccupied vertex and exactly one guard moves along one edge to the attacked
vertex.  The $\gamma$--$\theta$ conjecture asserts that
\(\gamma(G)=\gaminf(G)\) implies \(\gamma(G)=\theta(G)\).
We establish the following certificate-backed finite result:
no finite simple graph \(G\) on 12 vertices satisfies
\[
  \gamma(G)=\gaminf(G)=3<\theta(G).
\]
The result includes disconnected graphs.  The mathematical reduction sends
any putative graph to one of three overlapping cases in which
\(\overline G\) contains a hub-free induced \(C_5\), \(C_7\), or \(C_9\).
Each case is represented by a byte-bound propositional formula encoding the
graph, a nonempty closed family of dominating triples, and valid coloring
clauses implied by the obstruction to a three-clique partition.  Addition-only
reverse-unit-propagation refutations of all three formulas were replayed
independently.  We give the
coverage proof, exact formula and proof hashes, and reproduction commands.
This does not exclude parameter \(k\geq4\) at order 12, does not exclude
larger-order counterexamples, and does not resolve the universal conjecture.
\end{abstract}

\section{Introduction}

For a finite simple graph \(G\), let \(\gamma(G)\) be its domination number
and let \(\theta(G)\) be the minimum number of cliques in a partition of
\(\V(G)\).  In the standard one-guard eternal domination game, guards occupy
a dominating set, an attacker selects an unoccupied vertex, and the defender
moves one adjacent guard to the attacked vertex while retaining a dominating
configuration.  The least number of guards permitting defense against every
finite adaptive attack sequence is denoted by \(\gaminf(G)\).

Klostermeyer and MacGillivray stated the implication
\(\gamma=\gaminf\Rightarrow\gamma=\theta\), with a proof that was later
found to be incomplete \cite{KlostermeyerMacGillivray2009,
KlostermeyerMynhardt2015}.  Klostermeyer and Mynhardt consequently posed the
existence of a graph with
\[
  \gamma(G)=\gaminf(G)<\theta(G)
\tag{1.1}
\]
as an open problem \cite{KlostermeyerMynhardt2015,
KlostermeyerMynhardt2016}.  This is now commonly called the
\(\gamma\)--\(\theta\) conjecture.  MacGillivray, Mynhardt, and Virgile
reported an exhaustive search finding no counterexample through order 11
\cite{MacGillivrayMynhardtVirgile2022}.  Several proper graph classes are
also known, but those results do not settle (1.1) universally; see, for
example, the recent planar result \cite{Taletskii2024}.

Our main result is one precisely delimited slice at the next order.

\begin{theorem}\label{thm:main}
There is no finite simple graph \(G\) with \(|\V(G)|=12\) such that
\[
  \gamma(G)=\gaminf(G)=3<\theta(G).
\]
\end{theorem}

Theorem~\ref{thm:main} is not the assertion that there is no counterexample
of order 12: the common value in (1.1) could be at least four.  It is also
not an extension of the published all-parameter enumeration through order
11.  Its proof combines graph-theoretic reductions with exact finite
certificates.  Section~\ref{sec:reductions} gives the reductions,
Section~\ref{sec:split} proves that three complement templates are
exhaustive, Section~\ref{sec:encoding} gives the graph-to-formula coverage
argument, and Section~\ref{sec:certificates} records the checked
refutations.  Appendix~\ref{app:repro} gives the reproducibility record.

\section{Definitions and elementary reductions}\label{sec:reductions}

All graphs are finite, simple, and undirected.  For \(v\in\V(G)\), let
\(N(v)\) and \(N[v]\) denote its open and closed neighborhoods.  A set
\(D\subseteq\V(G)\) dominates \(G\) when
\(\bigcup_{v\in D}N[v]=\V(G)\).  Let \(\alpha(G)\) be the independence
number, and let \(i(G)\) be the minimum cardinality of a maximal independent
set.  The identity
\[
  \theta(G)=\chi(\overline G)
\tag{2.1}
\]
will be used throughout.

\begin{definition}
An \emph{eternal dominating family of size \(k\)} is a nonempty family
\(\mathcal D\subseteq\binom{\V(G)}k\) such that every
\(D\in\mathcal D\) dominates \(G\), and for every \(D\in\mathcal D\) and
every attacked vertex \(r\notin D\), there is some
\(u\in D\cap N(r)\) for which
\[
  (D-\{u\})\cup\{r\}\in\mathcal D.
\tag{2.2}
\]
The minimum such \(k\) is \(\gaminf(G)\).
\end{definition}

Thus attacks at occupied vertices are not part of the model, and (2.2)
moves exactly one guard along one edge.  No all-guards-move parameter is used
in this paper.

\begin{lemma}[independent-set forcing]\label{lem:forcing}
Let \(\mathcal D\) be an eternal dominating family of \(k\)-sets, and let
\(I\) be independent.  From every \(D\in\mathcal D\), attacks at currently
unoccupied vertices of \(I\) force a member \(D'\in\mathcal D\) with
\[
  |D'\cap I|=\min\{|I|,k\}.
\]
In particular, if \(|I|=k\), then \(I\in\mathcal D\).
\end{lemma}

\begin{proof}
Attack a vertex \(r\in I-D\).  A responding guard cannot move from another
vertex of \(I\), since \(I\) is independent.  The response therefore
increases \(|D\cap I|\) by one.  Repeat until either all guards or all
vertices of \(I\) are occupied.
\end{proof}

\begin{proposition}[parameter chain]\label{prop:chain}
For every nonempty graph \(G\),
\[
  \gamma(G)\leq i(G)\leq\alpha(G)\leq\gaminf(G)\leq\theta(G).
\tag{2.3}
\]
Consequently, if \(\gamma(G)=\gaminf(G)=k\), then
\[
  \gamma(G)=i(G)=\alpha(G)=\gaminf(G)=k.
\tag{2.4}
\]
\end{proposition}

\begin{proof}
Every maximal independent set dominates, and a maximum independent set is
maximal, proving the first two inequalities.  If an eternal family has
\(k<\alpha(G)\) guards, Lemma~\ref{lem:forcing} places all guards in a
maximum independent set while leaving another vertex of that set
unoccupied; the next attack there has no adjacent guard.  Hence
\(\alpha\leq\gaminf\).

For a partition of \(\V(G)\) into \(t\) cliques, place one guard in each
part.  An attack is answered by moving the guard within the attacked
vertex's clique, so \(\gaminf\leq t\).  Minimizing gives the final
inequality.  Equality (2.4) follows by squeezing.
\end{proof}

We need three further standard consequences.  Proofs are included because
they determine the exact search universe.

\begin{lemma}[induced-subgraph monotonicity]\label{lem:induced}
If \(J\) is a nonempty induced subgraph of \(G\), then
\(\gaminf(J)\leq\gaminf(G)\).
\end{lemma}

\begin{proof}
Let \(\mathcal D\) be an eternal family in \(G\), and among its configurations
retain those maximizing the number \(m\) of guards in \(\V(J)\).  Project
these configurations to \(\V(J)\).  For an attack inside \(J\), a response
moving a guard from outside \(J\) would increase the number of guards in
\(\V(J)\), contradicting maximality.  Hence every response stays inside
\(J\).  Because \(J\) is induced, it uses an edge of \(J\), preserves the
maximum \(m\), and projects to the required one-guard successor.

The same response argument shows that every projected state dominates
\(J\): each unoccupied vertex of \(J\) has the responding guard as a
neighbor in the projected state.  If \(m=0\), an attack at any vertex of
the nonempty graph \(J\) would necessarily move a guard into \(J\), again
contradicting maximality.  Thus \(m\geq1\), and the projected family is a
nonempty eternal family of \(m\)-sets.
\end{proof}

\begin{lemma}[component additivity]\label{lem:additivity}
If \(G_1,\ldots,G_t\) are the nonempty components of \(G\), then
\[
 \gamma(G)=\sum_j\gamma(G_j),\qquad
 \gaminf(G)=\sum_j\gaminf(G_j),\qquad
 \theta(G)=\sum_j\theta(G_j).
\tag{2.5}
\]
\end{lemma}

\begin{proof}
The assertions for \(\gamma\) and \(\theta\) follow by restriction to
components and union of componentwise solutions.  Products of eternal
families give the upper bound for \(\gaminf\).

For the reverse bound, choose one guard-count vector
\((|D\cap\V(G_1)|,\ldots,|D\cap\V(G_t)|)\) occurring in an eternal family,
and retain only configurations with that vector.  A legal move cannot cross
components, so this subfamily is closed under all attacks.  Its projection
to \(G_j\) is an eternal family with \(|D\cap\V(G_j)|\) guards.  Summing the
componentwise lower bounds proves equality.
\end{proof}

\begin{lemma}\label{lem:odd-antihole}
For odd \(n\geq5\),
\[
  \gaminf(\overline{C_n})=3.
\tag{2.6}
\]
\end{lemma}

\begin{proof}
The upper bound follows from
\(\theta(\overline{C_n})=\chi(C_n)=3\).  For the lower bound, label \(C_n\)
cyclically by \(0,\ldots,n-1\).  A pair fails to dominate
\(\overline{C_n}\) exactly when its two vertices have cyclic distance two.
If an eternal two-family existed, Lemma~\ref{lem:forcing} would put
\(\{0,1\}\) in it.  From \(\{0,d\}\), for odd
\(1\leq d\leq n-6\), attack \(d+2\).  Moving the guard at \(0\) leaves the
nondominating pair \(\{d,d+2\}\), so closure forces
\(\{0,d+2\}\).  Iteration reaches \(\{0,n-4\}\).  An attack at \(n-2\)
leaves a pair at cyclic distance two whichever guard moves, a
contradiction.  For \(n=5\), this final attack applies to the initial pair.
\end{proof}

This value is classical \cite{GoddardHedetniemiHedetniemi2005}; the argument
above fixes the convention used here.  We also need the following small
cycle value.

\begin{lemma}\label{lem:c7}
\(\gaminf(C_7)=4\).
\end{lemma}

\begin{proof}
A partition into three edges and one singleton gives the upper bound four.
Suppose an eternal family of dominating triples exists.  A dominating triple
has cyclic gap multiset \(\{2,1,1\}\) or \(\{2,2,0\}\).  Every triple of the
second type is equivalent to \(\{0,1,4\}\); attacking \(3\) forces the move
from \(4\), after which vertex \(5\) is undominated.  Thus no second-type
triple belongs to the family.

The maximum independent set \(\{0,2,4\}\) belongs to the family by
Lemma~\ref{lem:forcing}.  Attack \(1\).  Moving the guard at \(0\) leaves
vertex \(6\) undominated; moving the guard at \(2\) produces the forbidden
second-type triple \(\{0,1,4\}\).  This contradiction proves the lower
bound.
\end{proof}

\begin{lemma}[minimum counterexample parameter]\label{lem:min-k}
Every graph satisfying (1.1) has common value
\[
  \gamma(G)=\alpha(G)=\gaminf(G)\geq3.
\]
\end{lemma}

\begin{proof}
Equality follows from Proposition~\ref{prop:chain}.  Common value one makes
\(G\) complete.  Suppose the common value is two and
\(\theta(G)>2\).  For \(H=\overline G\),
\(\omega(H)=2<\chi(H)\).  The Strong Perfect Graph Theorem
\cite{ChudnovskyRobertsonSeymourThomas2006} gives an odd hole or odd
antihole in \(H\).  An odd hole in \(H\) gives an induced odd antihole in
\(G\), contradicting Lemmas~\ref{lem:induced} and
\ref{lem:odd-antihole}.  An odd antihole in \(H\) gives an odd hole in
\(G\).  At length at least seven this contradicts \(\alpha(G)=2\), and at
length five it again contradicts Lemma~\ref{lem:odd-antihole}, because
\(C_5\) is self-complementary.
\end{proof}

\section{The exhaustive three-template split}\label{sec:split}

Call a vertex outside an induced cycle a \emph{hub} if it is adjacent to
every vertex of the cycle.  A cycle is \emph{hub-free} if it has no hub.

\begin{proposition}\label{prop:connected}
If a graph \(G\) satisfies the hypothesis of Theorem~\ref{thm:main}, then
\(G\) is connected.
\end{proposition}

\begin{proof}
Let \(G_1,\ldots,G_t\) be its components and put
\(a_j=\gamma(G_j)\), \(b_j=\gaminf(G_j)\), and
\(c_j=\theta(G_j)\).  By (2.3) and (2.5),
\(\sum a_j=\sum b_j=3<\sum c_j\).  Nonnegativity of
\(b_j-a_j\) forces \(a_j=b_j\) for every \(j\), while the strict inequality
gives a component \(G_q\) with \(b_q<c_q\).  This component is itself a
counterexample, so Lemma~\ref{lem:min-k} gives \(a_q\geq3\).  Since all
nonempty components have positive domination number and \(\sum a_j=3\),
\(G_q=G\).
\end{proof}

\begin{proposition}[odd-wheel obstruction]\label{prop:wheel}
If \(\gaminf(G)=3\), then \(\overline G\) contains no induced odd wheel.
\end{proposition}

\begin{proof}
If \(\overline G\) induces \(K_1\vee C_{2q+1}\), then \(G\) induces
\(K_1\mathbin{\dot\cup}\overline{C_{2q+1}}\).  Lemmas
\ref{lem:additivity} and \ref{lem:odd-antihole} give eternal domination
number four for that induced subgraph, contradicting
Lemma~\ref{lem:induced}.
\end{proof}

\begin{theorem}[template split]\label{thm:split}
Let \(G\) satisfy the hypothesis of Theorem~\ref{thm:main}, and put
\(H=\overline G\).  Then \(H\) contains a hub-free induced
\(C_\ell\) for some
\[
  \ell\in\{5,7,9\}.
\tag{3.1}
\]
\end{theorem}

\begin{proof}
Proposition~\ref{prop:chain} gives
\[
  \omega(H)=\alpha(G)=3<\chi(H)=\theta(G).
\tag{3.2}
\]
Thus \(H\) is imperfect, and the Strong Perfect Graph Theorem supplies an
induced odd hole or odd antihole.  If an odd antihole has length \(2q+1\),
then \(q=\omega(\overline{C_{2q+1}})\leq3\).  It therefore has length five
or seven.  The five-antihole is \(C_5\).  An induced
\(\overline{C_7}\) in \(H\) would induce \(C_7\) in \(G\), contradicting
Lemmas~\ref{lem:induced} and \ref{lem:c7}.  Hence \(H\) contains an odd
hole.

The hole is hub-free by Proposition~\ref{prop:wheel}.  Moreover, because
\(\gamma(G)=3\), no pair dominates \(G\).  Equivalently, every pair of
vertices of \(H\) has a common \(H\)-neighbor.  The endpoints of a rim edge
of an induced hole have no common neighbor on the rim, so every rim edge has
a common neighbor outside the hole.  If exactly one vertex were outside,
that vertex would be adjacent to both endpoints of every rim edge, hence
would be a hub.  At least two vertices are therefore outside the hole.
Since \(|\V(H)|=12\), its odd length is at most ten, proving (3.1).
\end{proof}

\section{Exact finite encoding and coverage}\label{sec:encoding}

We now describe the formulas used to exclude the three cases in
Theorem~\ref{thm:split}.  The description is part of the coverage proof:
unsatisfiability of a formula is useful only after every graph in the
claimed branch is shown to produce a satisfying assignment.

Fix \(\ell\in\{5,7,9\}\).  Label an induced cycle of \(H=\overline G\) by
\(0,\ldots,\ell-1\) cyclically.  The endpoints \(0,1\) have a common
\(H\)-neighbor outside the rim; label one such vertex \(\ell\).  The
remaining vertices receive arbitrary labels.  The template fixes the cycle,
the edges \(0\ell,1\ell\), and a clause forbidding each outside vertex from
being a hub.

All three formulas use the same 6,886 Boolean variables:
\[
\begin{array}{rcll}
66   &:& e_{uv} & (uv\in E(H)),\\
660  &:& w_{ab,c} & (c\text{ witnesses a common }H\text{-neighbor of }a,b),\\
220  &:& f_D & (D\in\binom{[12]}3\text{ is in the eternal family}),\\
5,940&:& m_{D,r,u} & (u\in D\text{ responds to }r\notin D).
\end{array}
\tag{4.1}
\]
The base clauses assert:

\begin{enumerate}
\item no four vertices form a clique in \(H\);
\item every vertex pair has a common \(H\)-neighbor, with each true witness
      implying its two incident \(H\)-edges;
\item the labeled induced, hub-free \(C_\ell\) template;
\item every proper cut has a crossing \(G\)-edge, encoded as a negative
      \(H\)-edge;
\item every selected triple dominates \(G\);
\item at least one triple is selected;
\item for every selected \(D\) and every \(r\notin D\), at least one
      \(m_{D,r,u}\) is true, and every true move requires
      \(ur\in E(G)\) and selects the successor
      \((D-\{u\})\cup\{r\}\); and
\item every triangle of \(H\) is selected.
\end{enumerate}

The last clauses are redundant but sound: a triangle of \(H\) is an
independent triple of \(G\), so it lies in every eternal three-family by
Lemma~\ref{lem:forcing}.  Notice also that a legal move is encoded by
\(\neg e_{ur}\), because the edge variables describe \(H\), not \(G\).
The successor's family variable invokes the same domination clauses, so
every resulting configuration dominates.

To encode \(\chi(H)>3\), let \(c:[12]\to\{0,1,2\}\) be a coloring proper on
the forced template edges, and append
\[
  C_c=\bigvee_{\substack{u<v\\c(u)=c(v)}}e_{uv}.
\tag{4.2}
\]
This clause is false exactly when \(c\) is a proper coloring of \(H\).
Color names are canonicalized by first use.  A complete bank for the
\(\ell\)-template has
\[
  \frac{(2^\ell-2)3^{11-\ell}}{6}
\tag{4.3}
\]
rows: first color the odd rim, force the distinguished common neighbor to
use the third color, color the remaining vertices freely, and divide by the
free action of the six color permutations.

The \(C_5\) formula has one further, sound symmetry breaker.  For
\(v\in\{6,\ldots,11\}\), form the core signature
\[
  s(v)=(e_{0v},e_{1v},\ldots,e_{5v}).
\]
The complete formula is invariant under every permutation of these six
vertices, provided graph, witness, family, and move variables are all
relabelled.  Hence every model has an orbit representative satisfying
\[
  s(6)\leq_{\mathrm{lex}}s(7)\leq_{\mathrm{lex}}\cdots
  \leq_{\mathrm{lex}}s(11).
\tag{4.4}
\]
Each adjacent comparison is encoded by 63 clauses, one for every possible
first differing coordinate and common prefix, for 315 clauses in total.
Thus (4.4) is an equisatisfiable orbit restriction, not a logical
consequence silently added to the formula.

\begin{lemma}[graph-to-formula coverage]\label{lem:coverage}
Suppose that \(G\) is a connected 12-vertex graph with
\[
  \gamma(G)=\alpha(G)=\gaminf(G)=3<\theta(G)
\tag{4.5}
\]
and that \(\overline G\) contains a hub-free induced \(C_\ell\), where
\(\ell\in\{5,7,9\}\).  Then the graph and an eternal three-family induce a
satisfying assignment of the exact formula for that branch in
Table~\ref{tab:certificates}.
\end{lemma}

\begin{proof}
Use the labeling above and set \(e_{uv}\) from the actual edges of
\(H=\overline G\).  Equation (4.5) gives \(\omega(H)=3\), so the no-\(K_4\)
clauses hold.  Since no pair dominates \(G\), select an actual common
\(H\)-neighbor for every pair and set the corresponding witness variable.
Connectedness and hub-freeness satisfy the cut and template clauses.

Set \(f_D\) true exactly on an actual nonempty eternal three-family.  For
every selected \(D\) and \(r\notin D\), select one guard promised by (2.2).
This satisfies the domination, response, legal-edge, and successor clauses.
Lemma~\ref{lem:forcing} satisfies the triangle-selection clauses.

Finally, \(\theta(G)=\chi(H)>3\), so \(H\) fails every three-coloring.  It
therefore satisfies every appended clause (4.2), whether the branch uses a
complete bank or only a recorded subset.  For \(\ell=5\), permute the six
undistinguished outer vertices, together with all auxiliary variables, to
obtain an orbit-equivalent assignment satisfying (4.4).
\end{proof}

\begin{remark}
The \(C_5\) and \(C_7\) formulas use complete coloring banks of respectively
3,645 and 1,701 clauses.  The \(C_9\) certificate arose after 170 coloring
clauses in a counterexample-guided run.  Completeness of that 170-row list is
neither claimed nor needed: a target with \(\chi(H)>3\) satisfies every
valid clause (4.2), including those 170.
\end{remark}

\section{Proof certificates and the finite theorem}\label{sec:certificates}

A clause is a reverse-unit-propagation (RUP) consequence of a formula if
unit propagation after adjoining the negation of that clause derives a
contradiction.  Appending a RUP consequence preserves satisfiability.
Therefore an addition-only sequence of RUP consequences ending in the empty
clause is a checkable unsatisfiability certificate
\cite{WetzlerHeuleHunt2014}.

Table~\ref{tab:certificates} binds each formula and its accepted
addition-only proof.  Hashes are SHA-256.  The \(C_5\) proof is in binary
DRAT format; the other two are in ASCII DRAT format.  In every case the
accepted replay used only RUP: the checker reported zero RAT lemmas.

\begin{table}[ht]
\centering
\small
\caption{Exact branch formulas and accepted addition-only RUP proofs.}
\label{tab:certificates}
\begin{tabular}{@{}lrrr@{}}
\toprule
 & \(C_5\) & \(C_7\) & \(C_9\)\\
\midrule
variables & 6,886 & 6,886 & 6,886\\
clauses & 23,968 & 21,718 & 20,200\\
literals & 192,169 & 148,551 & 117,841\\
CNF bytes & 754,323 & 621,864 & 530,053\\
proof additions & 247,981 & 284,317 & 4,705\\
proof bytes & 6,337,621 & 18,093,724 & 65,906\\
\bottomrule
\end{tabular}
\end{table}

\noindent The formula hashes, in \(C_5,C_7,C_9\) order, are
\begin{align*}
&\texttt{c6a0811c718ff8e9352f253e4ce225ce2826def9c3e4a9cd55f0b0152703d104},\\
&\texttt{6a011e685e58ef517f2ab8253ca40987bd7b742a470bedbacdc3a5e94fc995a7},\\
&\texttt{2845f242a094484a8d114e70ca1a8678dfcff79fadd56bd57813e25c2e49523d}.
\end{align*}
The corresponding proof hashes are
\begin{align*}
&\texttt{c6c24853e30073e66fb396441edb176a0160d062a8558e25fa18a955f33927c3},\\
&\texttt{e8052df40d3e0c39b945a8735889039daba55eacc351e1822828b3d94f7baae9},\\
&\texttt{24c5647d3a57f2de221fba96747c618575a3aba086c5e4bca17aade55ce7d4ab}.
\end{align*}

For \(C_5\), an independent parser verified that the retained proof is
exactly the addition subsequence of the raw binary proof, with 245,439
deletions removed, and that its last addition is the empty clause.  A fresh
warning-fatal, forward, RUP-only replay used 10,912,555 resolution steps and
ended in \texttt{s VERIFIED}.  For \(C_7\), two strict warning-free RUP-only
replays accepted the 284,317-addition proof.  For \(C_9\), a separately
written standard-library checker validated all 4,705 RUP additions and the
final empty clause; a pinned external checker independently agreed.
The exact paths, checker hash, acceptance records, and commands appear in
Appendix~\ref{app:repro}.

\begin{proof}[Proof of Theorem~\ref{thm:main}]
Suppose \(G\) satisfies the displayed hypothesis.  By
Proposition~\ref{prop:connected}, \(G\) is connected, and by
Proposition~\ref{prop:chain} it satisfies (4.5).
Theorem~\ref{thm:split} gives a hub-free induced \(C_5\), \(C_7\), or
\(C_9\) in \(\overline G\).  Lemma~\ref{lem:coverage} then gives a model of
the exact corresponding formula in Table~\ref{tab:certificates}.  The
accepted RUP refutation proves that formula unsatisfiable, a contradiction.
\end{proof}

\begin{remark}[scope]
The theorem certifies the complete \((n,k)=(12,3)\) counterexample slice,
including disconnected graphs.  It does not certify the \(n=12,k\geq4\)
slices, any order \(n\geq13\), absence of counterexamples through order 12,
or the universal conjecture.
\end{remark}

\section{Research and verification provenance}

This manuscript and the associated research artifacts were produced with
substantial assistance from autonomous AI systems.  No mathematical claim
is accepted on model output alone.  The structural lemmas were subjected to
separate adversarial mathematical reviews; the exact formulas were
independently reconstructed; and the decisive proof streams were parsed and
replayed by programs that do not import the search implementation's
transition or encoding core.  Two independent reviews audited the final
composition from the three branch certificates to
Theorem~\ref{thm:main}.  Human authors remain responsible for the submitted
version and for inspecting the archived evidence.

\section*{Data availability}

The archive accompanying this manuscript contains the formulas, proofs,
standalone auditors, acceptance records, and exact hashes described below.
A permanent public archive identifier is to be inserted before submission.

\appendix

\section{Reproducibility and exact artifact bindings}\label{app:repro}

All paths below are relative to the campaign archive root
\texttt{gamma\_theta\_eternal\_domination/}.  The claim-level acceptance
record is
\begin{center}
\texttt{results/order12\_k3\_exclusion\_acceptance.json}.
\end{center}
The theorem and its two independent coverage reviews are
\begin{center}
\begin{minipage}{0.94\linewidth}
\texttt{math/lemmas/order12\_k3\_exclusion.md}\\
\texttt{reviews/order12\_k3\_exclusion\_hostile\_review.md}\\
\texttt{reviews/order12\_k3\_exclusion\_second\_review.md}.
\end{minipage}
\end{center}

The exact formulas are:
\begin{center}
\begin{minipage}{0.94\linewidth}
\texttt{results/synthesis\_k3\_hole5\_signature\_package/instance.cnf}\\
\texttt{results/synthesis\_k3\_template\_bank\_packages/hole7/instance.cnf}\\
\texttt{certificates/synthesis\_k3\_hole9\_orphan\_000170\_recovery/}\\
\hspace*{1em}\texttt{source/orphan-attempt-000170/instance.cnf}.
\end{minipage}
\end{center}
The accepted proofs are:
\begin{center}
\begin{minipage}{0.94\linewidth}
\texttt{results/synthesis\_k3\_hole5\_signature\_seed0\_600s\_binary/}\\
\hspace*{1em}\texttt{proof.additions.bdrat}\\
\texttt{certificates/synthesis\_k3\_hole7\_full\_bank\_seed0\_addition\_only\_v2/}\\
\hspace*{1em}\texttt{proof/addition-only.rup.drat}\\
\texttt{certificates/synthesis\_k3\_hole9\_orphan\_000170\_recovery/}\\
\hspace*{1em}\texttt{proof/addition-only.rup.drat}.
\end{minipage}
\end{center}

The pinned \texttt{drat-trim} executable has SHA-256
\[
\texttt{31df522b8b2b71acd357723b0e826cf488826ed78ad9e3a7bcad241271812beb}.
\]
The \(C_5\) source/run-gate revision is
\texttt{6f3ef0a0970b7214c34018fe32ea1ceeb5764d17}; its 12 run artifacts were
frozen without payload changes at
\texttt{dff45f4239e4acabc461533a0a213beec18ec56d}.  The complete C-035
acceptance package was committed at
\texttt{36d8191a}.

From the campaign archive root, one fail-closed command performs the
claim-level replay:
\begin{verbatim}
python3 -I -B repro/c035/replay.py --mode full \
  --output /fresh/path/c035-full-replay.json
\end{verbatim}

The wrapper first validates 91 exact file bindings, accepted Git objects,
branch activations, theorem scope, and manifest rows.  It then creates
fresh private temporary directories and sequentially runs the two
independent \(C_5\) audits, the sealed \(C_7\) proof replay, and the sealed
\(C_9\) proof replay.  It accepts only exact warning-free status records and
rehashes every input afterward.  A successful result has status
\texttt{PASS\_FULL\_C035\_REPLAY},
\texttt{claim\_status=CERTIFIED-FINITE}, and
\texttt{proofs\_freshly\_replayed=true}.  Time, memory, disk, and live-load
gates reject rather than weaken a run.  The fast mode validates metadata
only and explicitly reports
\texttt{NO\_MATHEMATICAL\_\allowbreak CLAIM}.

\bibliographystyle{plain}
\bibliography{references}

\end{document}
